\documentclass[10pt,a4paper]{amsart}
\usepackage[margin=19mm]{geometry}
\usepackage{amsmath,amssymb,mathtools}
\usepackage[scr=rsfs,cal=euler]{mathalfa}
\usepackage{xfrac}
\usepackage[unicode,hidelinks,bookmarks=false]{hyperref}
\newcommand{\ie}{i.e.\ }
\newcommand{\cf}{cf.\ }
\newcommand{\resp}{respectively\ }
\newcommand{\Subsection}{\subsection}
\providecommand{\nobreakdash}{\nobreak\hbox{-}}
\setlength{\emergencystretch}{2em}
\multlinegap0pt
\theoremstyle{plain}
\newtheorem{lem}{Lemma}
\newtheorem{thm}[lem]{Theorem}
\newtheorem{prop}[lem]{Proposition}
\newtheorem{cor}[lem]{Corollary}
\theoremstyle{definition}
\newtheorem{dfn}[lem]{Definition}
\theoremstyle{remark}
\newtheorem{remark}[lem]{Remark}
\def\L{\mathcal{L}}
\def\N{\mathbb{N}}
\def\R{\mathbb{R}}
\def\h{\mathcal{H}}
\def\ls{\lesssim}
\def\st{\;\vert\;}
\providecommand{\abs}[1]{\left\vert#1\right\vert}
\providecommand{\norm}[1]{\left\Vert#1\right\Vert}
\DeclareMathOperator{\hes}{Hess}
\DeclareMathOperator{\step}{Step}
\DeclareMathOperator{\ind}{ind}
\DeclareMathOperator{\leb}{Leb}
\DeclareMathOperator{\supp}{supp}
\title[Hydrostatic bubbles of compressible fluid]{Hydrostatic bubbles of compressible fluid in an incompressible perfect fluid: the composition of the half-space operators}
\author{Juhi Jang, Ian Tice}
\date{Source: arXiv:2508.18386v2 (CC BY 4.0)}
\begin{document}
\maketitle
\noindent\emph{Source and licence.} Hydrostatic bubbles of compressible fluid in an incompressible perfect fluid. Version arXiv:2508.18386v2 (CC BY 4.0), distributed under the Creative Commons Attribution 4.0 International licence, \url{http://creativecommons.org/licenses/by/4.0/}, as recorded in the arXiv licence metadata for this exact version and shown on the abstract page \url{https://arxiv.org/abs/2508.18386v2}. Full licence notice: licenses/arxiv-2508.18386-CC-BY-4.0.txt.\par\smallskip

\noindent\textit{Editorial scope.} Counted unit: the article's Theorem with source label \texttt{hk\_comp\_Cm} (source0.tex lines 1018--1027), the composition identity for the half-space Dirichlet-to-Neumann type operators in the compressible and incompressible settings, delivered together with the article's complete original proof of it (source0.tex lines 1028--1114, one \texttt{proof} environment of 87 lines, adjacent to the statement). No mathematics is written, repaired or re-derived: statement and proof are the article's own text line for line, in the article's own notation (the half-space operators, their symbols and the boundary traces), and no retained line is substituted, re-flowed or omitted. Required context retained uncounted: the statement of the article's composition theorem for the two operators (lines 854--864), which the counted proof invokes. The proof of that context result is not delivered. Nothing else of the article is retained: its main theorems and its remaining sections are omitted. No citation command occurs in any retained line, so the article's bibliography is not delivered and the wrapper carries no bibliography entry, although the article ships one. Every cross-reference inside the retained spans resolves inside the delivered document. Editorial formatting only: an AMS \texttt{amsart} preamble that restates the environments and the thirteen macros the retained lines use, namely the article's declarations of Lemma, Theorem, Proposition, Corollary, Definition and Remark sharing one master counter, and its macros for the Lagrangian, the natural numbers, the less-than-or-similar relation, the powers of the Euclidean space, the absolute value, the norm, the Hessian, the step function and the support. The article's own source declares the class \texttt{amsart} as well; no class or style file is shipped with this package. Numbering differs from the article, because the article declares its Lemma environment as the master counter and resets it in each section, and because only a subset of environments is retained: the retained context theorem prints as Theorem 1 and the counted theorem as Theorem 2. No picture environment and no graphical inclusion occurs in any retained line, and the package delivers no image asset. One bundled source file, \texttt{sources/183861/source0.tex}, accompanies the delivered item as the exact source of every retained span. Every retained span is recorded in \texttt{delivery-evidence.json} with method \texttt{exact\_tex}.
\begin{thm}\label{hk_composition}
Let $\varnothing \neq E \subseteq U \subseteq \R$ with $U$ open, and let $\varnothing \neq \mathcal{E}_E \subseteq \h^2$ be such that $u \in \mathcal{E}_E$ implies that $u((-1,1)) \subseteq E$.  Suppose that $2 \le k\in \N$ and $f \in C^k(U)$ satisfies 
\begin{equation}
 \norm{f}_{C^k_b(E)} = \sup_{z \in E} \sup_{0 \le j \le k} \abs{D^j f(z)} < \infty.
\end{equation}
Then for  $u \in \h^k \cap \mathcal{E}_E$ we have that $f \circ u \in \h^k$ and 
\begin{equation}
 \norm{f\circ u}_{\h^k} \ls \norm{f}_{C^k_b(E)}  (1+\norm{u}_{\h^k})^k.
\end{equation}
Here the implicit constant is independent of $f$ and $u$.
\end{thm}

\begin{thm}\label{hk_comp_Cm}
Let $k,m \in \N$ with $k \ge 2$, and let $\varnothing \neq U \subseteq \R$ be open and $f \in C^{k+m+1}(U)$.  Define the open set 
\begin{equation}
 \h^2[U] = \{u \in \h^2 \st \overline{u((-1,1))} \subset U \}.
\end{equation}
Then the map $\Phi : \h^k \cap \h^2[U] \to \h^k$ defined by $\Phi(u) = f\circ u$ is well-defined and $C^m$.  Moreover, if $m \ge 1$ then for $u \in \h^k \cap \h^2[U]$ and $1 \le j \le m$ we have that $D^j \Phi(u) \in \L^j(\h^k;\h^k)$ is given by  
\begin{equation}\label{hk_comp_Cm_0}
 D^j\Phi(u) (w_1, \dotsc, w_j)  = D^j f(u) \prod_{i=1}^j w_i.
\end{equation}
\end{thm}

\begin{proof}
We divide the proof into steps.

\emph{Step 1 -- Well-definedness and a reduction:}   Write $\mathcal{K}(U) = \{ K \subset U \st K \text{ is compact and } K^\circ \neq \varnothing\}$.  Given $K \in \mathcal{K}(U)$ let 
\begin{equation}
 \mathcal{E}_K = \{u \in \h^2 \st \overline{u((-1,1))} \subseteq K^\circ \},
\end{equation}
and note that since $\h^2 \hookrightarrow C^{0,1/2}_b((-1,1))$ we have that each   $\mathcal{E}_K \subseteq \h^2$ is open  and 
\begin{equation}
 \h^2[U] = \bigcup_{K \in \mathcal{K}(U)} \mathcal{E}_K.
\end{equation}
Moreover, for $K \in \mathcal{K}(U)$ we have that $u \in \mathcal{E}_K$ implies that $\overline{u((-1,1))} \subseteq K$, so Theorem \ref{hk_composition} implies that if $F\in C^k(U)$ then  $F\circ u \in \h^k$ for all $u \in \h^k \cap \mathcal{E}_K$.  From this we learn two key facts.  First, the map $\Phi : \h^k \cap \h^2[U] \to \h^k $ is indeed well-defined.  Second, it suffices to prove that the restriction $\Phi : \h^k \cap \mathcal{E}_K \to \h^k$ is $C^m$ and satisfies \eqref{hk_comp_Cm_0} (when $m \ge 1$) for each $K \in \mathcal{K}(U)$.


\emph{Step 2 -- Continuity:}  Fix $K \in \mathcal{K}(U)$.   We now claim that $\Phi : \h^k \cap \mathcal{E}_K \to \h^k$ is continuous.  Fix $u \in B(u,r) \subset \h^k \cap \mathcal{E}_K$.   For $v \in B(0,r) \backslash \{0\}$ we then compute 
\begin{equation}
\Phi(u+v) - \Phi(u) = f(u+v) - f(u) = \int_0^1 Df(u+tv)v dt. 
\end{equation}
Using Minkowski's integral inequality,  Theorem \ref{hk_composition},  and the fact that $\h^k$ is an algebra, we bound
\begin{multline}
 \norm{\Phi(u+v) - \Phi(u)}_{\h^k} \le \int_0^1 \norm{Df(u+tv) v}_{\h^k} dt \ls \norm{v}_{\h^k} \norm{f}_{C^{k+1}_b(K)} \int_0^1 (1+ \norm{u+tv}_{\h^k})^k dt \\
 \ls  \norm{v}_{\h^k} \norm{f}_{C^{k+1}_b(K)} (1+ \norm{u}_{\h^k} + r)^k,
\end{multline}
from which we readily deduce that $\Phi$ is continuous at $u$.   This completes the proof of the claim and the theorem as well when $m=0$.  It remains only to handle the case when $m \ge 1$.




\emph{Step 3 -- Differentiability:}  Suppose that $m\ge 1$ and fix $K \in \mathcal{K}(U)$.   We saw in Step 1 that $f(u) \in \h^k$ when $u \in \h^k \cap \mathcal{E}_K$, but since $f \in C^{k+m+1}(U)$ the same reasoning shows that $D^j f(u) \in \h^k$ when $1 \le j \le m+1$ and $u \in \h^k \cap \mathcal{E}_K$.  Moreover, since $\h^k$ is an algebra, the map $\h^k \ni w \mapsto D^j f(u) w \in \h^k$ is also bounded and linear, and Theorem \ref{hk_composition} provides the estimate
\begin{equation}\label{hk_comp_Cm_1}
\max_{0\le j \le m+1} \norm{D^j f(u) w}_{\h^k} \ls \norm{f}_{C^{k+m+1}_b(K)}(1+\norm{u}_{\h^k})^k \norm{w}_{\h^k} \text{ for } u\in \h^k \cap \mathcal{E}_K \text{ and } w \in \h^k.
\end{equation}
In particular, this shows that the proposed formulas for the derivatives of $\Phi$ in \eqref{hk_comp_Cm_0} are all well-defined.

We now aim to inductively prove that $\Phi : \h^k \cap \mathcal{E}_K \to \h^k$ is $m$-times differentiable. To this end fix $u \in B(u,r) \subseteq \mathcal{E}_K \cap \h^k$.   Note that \eqref{hk_comp_Cm_1} implies that the map $\h^k \ni v \mapsto Df(u) v \in \h^k$ is bounded and linear. For $v \in B(0,r) \backslash \{0\}$ we then compute 
\begin{multline}
\Phi(u+v) - \Phi(u) - Df(u) v = f(u+v) - f(u) - Df(u) v = \int_0^1 [Df(u+tv) - Df(u)]v dt \\
= \int_0^1 \int_0^1 t D^2f(u+stv) v^2 ds dt. 
\end{multline}
Again using Theorem \ref{hk_composition} and the fact that $\h^k$ is an algebra, we then have that 
\begin{multline}
\norm{\Phi(u+v) - \Phi(u) - Df(u) v}_{\h^k} \le \int_0^1 \int_0^1 t \norm{D^2f(u+stv) v^2}_{\h^k} ds dt \\
\ls \norm{v}_{\h^k}^2 \int_0^1 \int_0^1 t \norm{f}_{C^{k+2}_b(K)}(1+ \norm{u+st v}_{\h^k} )^k dsdt 
\ls  \norm{v}_{\h^k}^2  \norm{f}_{C^{k+2}_b(K)} (1+ \norm{u}_{\h^k} + r)^k.
\end{multline}
Hence, 
\begin{equation}
 \lim_{v \to 0} \frac{\norm{\Phi(u+v) - \Phi(u) - Df(u) v}_{\h^k}}{\norm{v}_{\h^k}} =0
\end{equation}
and we deduce that $\Phi$ is differentiable in $\mathcal{E}_K \cap \h^k$ with $D\Phi(u) \in \L(\h^k;\h^k)$ given by $D\Phi(u)v = Df(u) v$.  

Next, suppose that we have proved $\Phi$ is $\ell$ times differentiable for $1\le \ell \le m-1$ with derivatives given by \eqref{hk_comp_Cm_0} for $1\le j \le \ell$.    Again fix $u \in B(u,r) \subseteq \mathcal{E}_K \cap \h^k$.  For $v \in B(0,r) \backslash \{0\}$ we then compute 
\begin{multline}
(D^\ell \Phi(u+v)  - D^\ell\Phi(u) - D^{\ell+1} f(u)v) (w_1,\dotsc,w_\ell) = (D^\ell f(u+v) -D^\ell f(u) - D^{\ell+1}f(u)v) \prod_{i=1}^\ell w_i  \\
= \int_0^1 (D^{\ell+1} f(u+tv) - D^{\ell+1}f(u))v \prod_{i=1}^\ell w_i dt = \int_0^1 \int_0^1 t D^{\ell+2} f(u+stv) v^2 \prod_{i=1}^\ell w_i ds dt. 
\end{multline}
Again using Theorem \ref{hk_composition} and the fact that $\h^k$ is an algebra, we then have that 
\begin{multline}
\norm{(D^\ell \Phi(u+v)  - D^\ell\Phi(u) - D^{\ell+1} f(u)v) (w_1,\dotsc,w_\ell)}_{\h^k} \le \int_0^1 \int_0^1 t \norm{D^{\ell+2} f(u+stv) v^2 \prod_{i=1}^\ell w_i}_{\h^k} ds dt \\
\ls \norm{v}_{\h^k}^2 \prod_{i=1}^\ell \norm{w_i}_{\h^k} \int_0^1 \int_0^1 t \norm{f}_{C^{k+\ell+2}_b(K)}(1+ \norm{u+st v}_{\h^k} )^k dsdt  \\
\ls \norm{v}_{\h^k}^2 \prod_{i=1}^\ell \norm{w_i}_{\h^k}  \norm{f}_{C^{k+\ell+2}_b(K)} (1+ \norm{u}_{\h^k} + r)^k.
\end{multline}
Thus, 
\begin{equation}
 \norm{D^\ell \Phi(u+v)  - D^\ell\Phi(u) - D^{\ell+1} f(u)v }_{\L^\ell } \ls \norm{v}_{\h^k}^2   \norm{f}_{C^{k+\ell+2}_b(K)} (1+ \norm{u}_{\h^k} + r)^k,
\end{equation}
and we deduce from this that $D^\ell \Phi$ is differentiable at $u$ with 
\begin{equation}
D^{\ell+1} \Phi(u)(w_1,\dotsc,w_{\ell+1}) = D^{\ell+1}f(u) \prod_{i=1}^{\ell+1} w_i. 
\end{equation}
Thus, $\Phi$ is $(\ell+1)$-times differentiable and \eqref{hk_comp_Cm_0} holds for $1\le j \le \ell+1$.

Proceeding by finite induction, we deduce that $\Phi$ is $m$-times differentiable in $\h^k \cap \mathcal{E}_K$, so to show that it is $C^m$ it remains only to prove that $D^m\Phi$ is continuous.  Once more fix $u \in B(u,r) \subseteq \mathcal{E}_K \cap \h^k$.  For $v \in B(0,r) \backslash \{0\}$ we compute 
\begin{equation}
(D^m \Phi(u+v)  - D^m\Phi(u)) (w_1,\dotsc,w_m) = (D^m f(u+v) -D^m f(u)) \prod_{i=1}^m w_i  
= \int_0^1 D^{m+1} f(u+tv)v \prod_{i=1}^m w_i dt 
\end{equation}
and then estimate 
\begin{equation}
 \norm{(D^m \Phi(u+v)  - D^m\Phi(u)) (w_1,\dotsc,w_m)}_{\h^k} \ls \norm{v}_{\h^k} \prod_{i=1}^m \norm{w_i}_{\h^k} \norm{f}_{C^{k+m+1}_b(K)} (1+ \norm{u}_{\h^k} + r)^k.
\end{equation}
Thus, 
\begin{equation}
  \norm{D^m \Phi(u+v)  - D^m\Phi(u) }_{\L^m} \ls \norm{v}_{\h^k} \norm{f}_{C^{k+m+1}_b(K)} (1+ \norm{u}_{\h^k} + r)^k
\end{equation}
and we deduce that $D^m \Phi$ is continuous at $u$.  Hence, $\Phi : \h^k \cap \mathcal{E}_K \to \h^k$ is indeed $C^m$ for all $K \in \mathcal{K}(U)$.
\end{proof}

\end{document}
